% =====================================================================
\section{The Walkable Bunker}\label{sec:walk}
The command bunker is underground, so the player walks into it. Two CPU-only, unit-tested pieces implement this (T50--T54, T60--T69, T120--T126): the \code{WalkSpace}
(floors, walls and indoor volumes, \code{scene/WalkSpace.h}, built by \code{Bunker.cpp}) and the \emph{walk camera} (\code{Camera::moveWalk}).

\subsection{WalkSpace}
\begin{itemize}
  \item \textbf{Floors.} A \code{WalkFloor} is a rectangle with a floor height at its low and high edge along a slope axis (flat, $x$ or $z$) and a ceiling
  height. The bunker defines 12: the entrance ramp (a linear slope along $z$ from the surface at $z=13$ to $-3.4$\,m at $z=6$), a level
  approach, the door zone, the tunnel, the CIC opening, the command room ($x\in[-5.2,5.2]$, $z\in[-8.2,-2.1]$), the bunk room and doorway, the generator room and doorway, the shaft entrance and the exit shaft.
  Outside every floor the walker stands on the terrain.
  \item \textbf{Which floor.} \code{floorHeight(x,z,feetY)} returns the \emph{highest} floor containing $(x,z)$ that is at most a 0.45\,m step above the feet
  and whose ceiling is above the feet, so a walker on the mound above the rooms stays on the terrain and one in the tunnel is not pulled onto the roof.
  \item \textbf{Walls.} 36 axis-aligned boxes (retaining and cross walls, furniture, jambs). \code{collide()} pushes a vertical cylinder (radius 0.3\,m) out of every box, in two passes
  to settle corners; boxes below the feet or above the head are ignored.
  \item \textbf{Dynamic obstacles.} The closed blast door and a lowered barrier arm are extra boxes each frame (T64, T121).
  \item \textbf{Indoor volumes.} Five boxes (tunnel, CIC, bunk, generator, shaft) form the shader's indoor mask (Section~\ref{sec:shading}).
\end{itemize}
The bunker part budget in test T66 reads: \bunkerBudget. T69: 12 floors, 36 walls, 5 indoor volumes.

\subsection{Walk camera}
Key \texttt{C} cycles Orbit, Free, Chase, UAV CAM and Walk (T55). In walk mode the eye is 1.7\,m above the feet (T51: maximum deviation 0.013\,m over 2\,400 steps), the field of
view $60^\circ$, the speed 1.6\,m/s (4.5 with Shift). Steps are split into $\le0.15$\,m pieces so a fast frame cannot tunnel through a door, and the feet follow the floor after each piece
(T120: ramp deviation 0\,m). The eye stays 0.15\,m below the ceiling (T123), pitch is clamped to $\pm80^\circ$ (T54), the position to a 52\,m radius (T53), and the walker cannot fall into the trench (T52).
A small head bob is added to the rendered eye only.

\subsection{Verified walk-through}
T62 walks from the ramp top to the CIC centre in 367 steps with 0 wall contacts, a maximum step of 0.023\,m and a minimum head room of 2.60\,m. T121: with the doors closed the walker stops at
$z=4.4$\,m (door plane $z=4$) and stays on the ramp ($|x|\le1.2$\,m). T122: with the doors open he reaches the CIC ($z=-6.0$\,m, feet $-3.40$\,m, smallest ceiling gap 0.90\,m).
T126: the director's bunker shot ends in the rooms with a maximum vertical step of 0.118\,m, and T129 checks that the director opens the doors early.

\begin{figure}[h]
  \centering
  \begin{subfigure}{0.49\linewidth}\picfour{bunker_ramp}{\linewidth}\caption{Entrance ramp (walk camera)}\end{subfigure}
  \begin{subfigure}{0.49\linewidth}\picfour{bunker_door}{\linewidth}\caption{Blast doors}\end{subfigure}\\[1mm]
  \begin{subfigure}{0.49\linewidth}\picfour{bunker_tunnel}{\linewidth}\caption{Tunnel}\end{subfigure}
  \begin{subfigure}{0.49\linewidth}\picfour{bunker_cic}{\linewidth}\caption{Command room (CIC)}\end{subfigure}
  \caption{Walking into the bunker (captures \code{bunker\_*}; eye heights 1.7\,m).}
  \label{fig:bunkerwalk}
\end{figure}
